% Transcription — fonds Alexandre Grothendieck, shelfmark 15, batch 12 (pages 221-227).
% Established from the facsimile archives/batches/15/batch-12.pdf (a mirror of
% https://grothendieck.umontpellier.fr/15.pdf), per the skill
% transcribe-grothendieck. The batch file carries no cover sheet: PDF page N
% is page 220+N of the folder (the Montpellier footer prints « 221 » ... « 227 »),
% and the archivists' pencil numbers bottom-left agree leaf by leaf (221 to
% 227 all legible). This is the last batch of the folder: seven pages, not
% twenty.
%
% Pass: Opus 5.5 (claude-opus-5-5), 2026-09-23 — first pass, unchecked against the pages by a human.
%
% Three of the seven pages are transcribed. Absent: pages 221, 223, 225 and
% 227, leaves of an unrelated typescript — a Bourbaki draft marked « n°370 »
% (Haar measure: relatively invariant measures, the modulus, Theorem
% 1 on existence and uniqueness of Haar measure), scanned upside down,
% nothing in his hand. They are not the folder's own « tapuscrit ».
%
% What the batch is. Three fast handwritten leaves, not paginated by him and
% without headings, on the fibre functors of motives over a number field and
% the Hodge realisation:
%   222    (continuing presumably from batch 11 — the page opens mid-thought,
%          « Il s'agirait de construire ») X projective smooth over k with
%          H^*(X_v, Z) and H^*(X_v', Z) not isomorphic as rings, « mettre Serre
%          à contribution »; margin: can Z be replaced by R? Then: k a number
%          field, v_0 a place, M a motive over k, G_{M,v_0} an algebraic group
%          over Z; every other place v gives a « fibré principal homogène » P_v
%          under G_{v_0}; conjectured « plus haut » (M general enough):
%          v ≠ v' ⟹ P_v ≠ P_v'. G_{M,v_0}(Z) = π^1_M(k,v), a discrete group;
%          « Vers : groupe fondamentaux (…) discrets ». Rotated margin: the
%          condition of an isomorphism of functors ξ(M)⊗_Z C ≃ M_∞⊗_k C
%          compatible with Riemann–Hodge conditions may be a serious
%          restriction.
%   224    A motive (?) is far from determined by its fibre functor into
%          Z-modules; what if one takes the functor with its Hodge structure?
%          (interlinear: true for … kählérien). A block reducing to k_0 = the
%          algebraic closure of Q in k, X of finite type over Q, then over C,
%          then to P^r_C, is boxed and cancelled by diagonal strokes. Then:
%          reduction to X = P^1_C − (0) − (1) − (∞), with the motive of a
%          family of elliptic curves over X whose invariant is the
%          indeterminate (?). Margin: a small diagram of fields Q, k_0, k.
%   226    Which multiplicative functors from motives (over C) to Hodge
%          structures are … (margin: « en vertu de la conjecture de Hodge »);
%          can a multiplicative functor be taken arbitrary « ?!? »; link with
%          the construction of moduli schemes for polarised motives; Y → Modules;
%          « Kif kif »; « Moralement … si Y est le spectre d'un corps ….. »,
%          where the page, the batch and the folder stop.
%
% Notation fixed for the batch. k the number field, v_0, v, v' its places,
% G_{M,v_0} (he writes G_{Mv_0}, the subscript run together) and G_{v_0};
% P_v the torsor; \pi^1_M(k,v) with a bold M, as written; \xi(M) and
% M_\infty on page 222; \mathbb{P}^1_{\mathbb{C}} on page 224. Plain square
% brackets in the text are his interlinear additions (as in folder 14); the
% editorial additions are \add{}.
%
% Legibility. This is the fast hand of a drafting session: the mathematical
% skeleton carries, the connective prose often does not, and the cancelled
% block on page 224 is only partly read. Nothing on these leaves is dated.
%
% The dating is the inventory's own for the folder, copied verbatim. Its
% group, [10-18] « Théorie des motifs », is dated [à partir de 1958-à partir
% de 1982].
\documentclass[11pt,a4paper]{article}
% Le préambule commun à toutes les transcriptions du fonds Grothendieck.
%
% Il tient en une idée : **l'incertitude se compose, elle ne se résout pas.**
% Une transcription qui lisse un mot douteux a détruit la seule chose pour
% laquelle on la faisait — un lecteur doit pouvoir savoir, à chaque endroit,
% ce qui était sur le feuillet et ce qui a été ajouté. Les six macros qui
% suivent sont donc l'appareil critique, et non de la décoration.
%
% Les mêmes commandes sont reconnues par `scripts/render.mjs`, qui produit la
% vue de lecture du site. Une macro ajoutée ici doit l'être aussi là-bas,
% faute de quoi le rendu échouera bruyamment — ce qui est voulu : un rendu
% silencieusement faux est pire qu'un rendu absent.

\ProvidesPackage{grothendieck}[2026/08/08 Transcriptions du fonds Grothendieck]

\usepackage{amsmath,amssymb,amsthm}
\usepackage{mathrsfs}
\usepackage{stmaryrd}
% Les diagrammes commutatifs sont partout dans ce fonds : c'est la langue
% même de la géométrie algébrique de Grothendieck, et une transcription qui
% les rendrait en prose ne transcrirait rien.
\usepackage{tikz-cd}
% Français par défaut — c'est la langue des feuillets — mais l'anglais est
% chargé aussi : la traduction et le résumé sont des documents anglais, et la
% typographie française (espaces avant les deux-points) y serait une faute.
% Ces documents basculent par \selectlanguage{english} en tête de corps.
\usepackage[english,french]{babel}
\usepackage{fontspec}
\usepackage{geometry}
\geometry{a4paper,margin=25mm,marginparwidth=28mm}
\usepackage{marginnote}
\usepackage{xcolor}
% ulem, not soul. `soul`'s \st and \ul are built on a character-by-character
% scan that XeLaTeX's font machinery breaks: the document fails with an
% undefined control sequence rather than degrading. ulem does the same two jobs
% and is Unicode-engine safe. `normalem` stops it hijacking \emph, which the
% transcriptions use for Grothendieck's own emphasis.
\usepackage[normalem]{ulem}
\usepackage{hyperref}
\hypersetup{colorlinks=true,linkcolor=black,urlcolor=black,pdfborder={0 0 0}}

% --- Métadonnées du lot -----------------------------------------------------
% Reprises telles quelles par le script de rendu, qui les affiche en tête de
% la vue de lecture. Elles fixent ce que le fichier prétend couvrir : sans
% elles, une transcription flotte sans point d'attache dans un fonds de seize
% mille feuillets.

\newcommand{\folder}[1]{\gdef\grothFolder{#1}}
\newcommand{\batch}[1]{\gdef\grothBatch{#1}}
\newcommand{\leaves}[2]{\gdef\grothFirst{#1}\gdef\grothLast{#2}}
\newcommand{\foldertitle}[1]{\gdef\grothFoldertitle{#1}}
\gdef\grothFolder{?}\gdef\grothBatch{?}\gdef\grothFirst{?}\gdef\grothLast{?}\gdef\grothFoldertitle{}

% --- Le résumé --------------------------------------------------------------
%
% Il ouvre la lecture modernisée, sous le titre « Résumé », et il est composé
% en petit. Ce n'est pas un ornement : c'est le seul passage du document qui
% ne soit pas des mathématiques, et un lecteur doit pouvoir le reconnaître
% avant de l'avoir lu — puis le sauter s'il connaît déjà le sujet, ou s'y
% arrêter s'il ne le connaît pas. Le filet qui le suit marque la reprise du
% corps.

\newenvironment{resume}
  {\small\setlength{\parskip}{0.45em}}
  {\par\medskip\hrule\medskip}

% --- Mots-clés ---------------------------------------------------------------
%
% En anglais, en clôture du résumé de la lecture modernisée : le vocabulaire
% moderne sous lequel chercher ce que les feuillets construisent. Le manifeste
% du site (`npm run manifest`) les extrait pour étiqueter la cote entière —
% cette ligne est la source unique des tags, il n'y a pas de fichier à côté.
\newcommand{\keywords}[1]{\par\smallskip\noindent
  {\footnotesize\textsc{Keywords} --- \textit{#1}}\par}

% --- L'appareil critique ----------------------------------------------------

% Le numéro de feuillet, en marge. C'est l'ancre qui permet de régler un
% désaccord sur une formule en regardant *un* feuillet plutôt qu'une chemise
% de six cents. Le site s'en sert aussi pour tourner les pages du fac-similé
% quand on fait défiler la transcription.
\newcommand{\leaf}[1]{%
  \marginnote{\footnotesize\color{gray!70}\textsf{#1}}%
  \label{leaf:#1}\ignorespaces}

% Illisible : on ne devine pas. Un mot inventé qui se lit comme les autres est
% le pire résultat possible.
\newcommand{\ill}{\textcolor{red!60!black}{[\,\ldots\,]}}

% Lecture douteuse : le mot est donné, et signalé comme incertain.
\newcommand{\uncertain}[1]{\uline{#1}}

% Ajout de l'éditeur — une lettre manquante, un mot sous-entendu.
\newcommand{\add}[1]{\textcolor{blue!50!black}{[#1]}}

% Biffé par Grothendieck lui-même. Ce qu'il a rayé fait partie du document :
% une rature dit souvent où la pensée a changé de direction.
\newcommand{\struck}[1]{\sout{#1}}

% Note du transcripteur, sur l'état matériel du feuillet.
\newcommand{\note}[1]{\par\smallskip{\small\color{gray!60!black}\textit{[#1]}}\par\smallskip}

% Note marginale de Grothendieck, distincte du corps du texte.
\newcommand{\marginal}[1]{\par\smallskip\noindent\hspace{1em}%
  {\small\color{gray!60!black}#1}\par\smallskip}

% --- Titre ------------------------------------------------------------------

\AtBeginDocument{%
  \begin{center}
    {\large\bfseries Fonds Alexandre Grothendieck --- cote n\textsuperscript{o}~\grothFolder}\\[2pt]
    {\itshape\grothFoldertitle}\\[4pt]
    {\small feuillets \grothFirst--\grothLast\ (lot \grothBatch)}\\[2pt]
    {\footnotesize\color{gray!60!black}Transcription établie d'après le fac-similé numérisé
     par l'Université de Montpellier.\\
     Les lectures incertaines sont soulignées, les illisibles notées [\,\ldots\,],
     les ajouts entre crochets.}
  \end{center}
  \vspace{1em}}

\endinput


\folder{15}
\batch{12}
\pages{221}{227}
\dating{1965-[vers 1977]}
\watermark{Édition de démonstration}
\foldertitle{Théorie arithmétique. Théorie de Galois des motifs (1965) : notes manuscites (s.d.), tapuscrit (s.d.).}

\begin{document}

\page{222}
\note{le feuillet commence au milieu d'un raisonnement ; la phrase qui
précède est vraisemblablement au lot 11 (le feuillet 221 est une page
dactylographiée étrangère)}

\marginal{Peut-on \uncertain{remplacer} $\mathbb{Z}$ par $\mathbb{R}$ ??}

Il \uncertain{s'agirait} \uncertain{de} construire $X$ projectif lisse sur
$k$, tel que $H^{*}(X_{v}, \mathbb{Z})$ et $H^{*}(X_{v'}, \mathbb{Z})$
non isomorphes comme \emph{anneaux}. On devrait mettre Serre à contribution.
\note{les indices de $X$ sont mal formés ; $v$, $v'$ sont lus d'après la
suite du feuillet, où $v$ désigne une place de $k$. Un double trait
horizontal sépare ce paragraphe du suivant}

Soit $k$ un corps de nombres, $v_{0}$ une place de $k$, $M$ un motif sur
$k$, d'où $G_{M,v_{0}} = G_{v_{0}}$ groupe algébrique \emph{défini sur}
$\mathbb{Z}$. Toute autre place \struck{\ill{}} $v$ définit un fibré
principal homogène $P_{v}$ sous $G_{v_{0}}$, \ill{} \ill{} \ill{} été
\struck{dit} [conjecturé] plus haut, [\uncertain{en admettant} que $M$ est
vraiment assez général,] $v \neq v' \Rightarrow P_{v} \neq P_{v'}$. On
\uncertain{peut} \uncertain{aller} plus loin et se demander
(\struck{en} $G_{M}$ assez grand) \struck{\ill{}} \uncertain{si} les
fibrés principaux homogènes $P_{v}$ sous $G_{v_{0}}$
d'une place $v$ \struck{\uncertain{réelle}} / \struck{les} \ill{}
\struck{\ill{}} \ill{} \ill{}.
\marginal{Attention, la condition d'avoir un isomorphisme de foncteurs
$\xi(M) \otimes_{\mathbb{Z}} \mathbb{C} \simeq M_{\infty} \otimes_{k}
\mathbb{C}$, fonctoriel, satisfaisant aux conditions de
\uncertain{compatibilité} genre Riemann--Hodge, risque d'être une condition
restrictive sérieuse…}
\note{cette note marginale, écrite en biais dans la marge gauche, est
appelée par un signe à tête ronde qui barre le mot biffé après « place $v$ » ;
la fin de la phrase est illisible}
\[
  G_{M,v_{0}}(\mathbb{Z}) = \pi^{1}_{\mathbf{M}}(k, v) \, .
\]
un groupe \uncertain{a priori} un \ill{} \ill{} \ill{}

\uncertain{Vers} : groupe \emph{fondamentaux}
[(\uncertain{arithmétiques})] discrets.
\note{« fondamentaux » est souligné ; le mot entre parenthèses est ajouté
au-dessus de la ligne et renvoyé par un trait avant « discrets »}

\page{224}
\marginal{\struck{\ill{}}}
\note{note de quatre ou cinq mots dans le coin supérieur gauche, biffée et
illisible}

On a vu plus haut qu'un \ill{} est très loin d'être caractérisé par le
foncteur fibre correspondant, \struck{Mais} considéré comme foncteur dans
les $\mathbb{Z}$-modules. Mais que se passe-t-il si on le regarde comme
foncteur dans les \ill{} \struck{\ill{}} \ill{} et structure de Hodge ?
\struck{\ill{}} [Vrai pour un \ill{} \uncertain{kählérien}, et
\uncertain{par là}]
\note{la phrase ajoutée dans l'interligne court au-dessus du début du
passage biffé qui suit ; sa liaison avec lui n'est pas claire}

\marginal{Attention, il faut faire \ill{} \ill{} \ill{} \ill{} \ill{} \ill{}
\ill{} \ill{} structure \ill{} de Hodge …}
\note{la note marginale est écrite en biais dans la marge gauche ; au-dessous,
un petit croquis relie $\mathbb{Q}$, $k_{0}$ (?), $k$ et $X$ par des
traits, sans doute la tour de corps du passage biffé ci-dessous}

\struck{\ill{}, on est ramené au cas où \ill{} \struck{plus} coïncide avec
\ill{}, $k_{0}$ = clôture alg.\ de $\mathbb{Q}$ dans $k$, et puis \ill{} au
\ill{} $X$ est de type fini sur $\mathbb{Q}$, puis on \ill{} de base \ill{}
\ill{} $X$ est connexe de type fini sur $\mathbb{C}$, [\uncertain{ordinaires}]
\ill{} de $X$ \ill{} \ill{} \ill{} quasi-projectif), \ill{} $X$ \ill{}
\ill{} au \ill{} $X = \mathbb{P}^{r}_{\mathbb{C}}$}
\note{tout ce passage est encadré et biffé de longs traits diagonaux ; les
lignes sont en outre barrées une à une. Seuls les fragments donnés se
lisent ; l'ordre des lignes est celui de la page}

Il semble \uncertain{qu'on se} ramène facilement au [\ill{} \ill{} \ill{}
\ill{} en \ill{},] cas de
\[
  X = \mathbb{P}^{1}_{\mathbb{C}} - (0) - (1) - (\infty) ,
\]
où la théorie \ill{} \struck{\uncertain{doit}} le donner --- [\ill{}
\uncertain{comme} motif celui défini par un schéma en courbes elliptiques
sur $X$, d'invariant \uncertain{l'indéterminée}].
\note{« facilement » est relié par un trait au bas du cadre biffé, là où se
lit $X = \mathbb{P}^{r}_{\mathbb{C}}$. L'exposant de $\mathbb{P}$ est
lu $1$, conformément aux trois points retirés}

\page{226}
\marginal{en vertu de la conjecture de Hodge}
\note{note encerclée dans le coin supérieur gauche, sans renvoi explicite}

D'autre part, \emph{quels} sont les foncteurs multiplicatifs de Motifs,
à valeurs dans [(i.e.\ pour les motifs sur $\mathbb{C}$)] les structures
de Hodge, qui sont \uncertain{motivés}, i.e.\ \uncertain{les} plus
\uncertain{ordinaires} ? Plus généralement, dans quelle mesure un
\ill{} [\ill{} $X - Y$ \uncertain{et il}] \ill{} \ill{} le foncteur sur les
motifs [(\uncertain{loc.\ cit.})], et dans quelle mesure un
\struck{foncteur} [multiplicatif] peut-il être pris arbitrairement ?!?
\note{« un morphisme » se lit peut-être au début de la ligne ; au-dessus,
« $X - Y$ » et une ligne illisible sont ajoutés dans l'interligne}

C'est lié au Pb de \emph{construction des schémas de modules} pour les
motifs [\ill{}] polarisés ; si $Y$ admet un \ill{} polarisé \ill{} tel que
$Y \to \mathrm{Modules}$ soit \ill{} (\ill{} …) une \uncertain{immersion})
\struck{et} alors on \ill{} \ill{}. Kif kif si $Y$ est \ill{} $\leftarrow$
\ill{} de \uncertain{problèmes} \ill{} \uncertain{ayant} \ill{} la \ill{}
propriété. Moralement \ill{} \uncertain{particulier} si $Y$ est le spectre
d'un corps …..
\note{la phrase s'arrête sur ces points de suspension, aux deux tiers du
feuillet ; c'est la dernière page manuscrite du dossier}

\end{document}
